\honest{Fake Causality}{Eliezer Yudkowsky}{2007}

Phlogiston was the eighteenth century's answer to what fire is and why wood becomes ash, ``and that was it.'' It explained flames, ash, flames going out in closed containers, and why charcoal leaves little residue. But no one used it to predict: ``You looked at the result first,'' and then explained it. ``It could explain everything.'' \nb{The historical record is more mixed. The theory forbade outcomes: burned metals should have lost phlogiston, and the weight some of them gained counted against it. In 1783 Priestley predicted, and then saw, that a metal calx heated in inflammable air, which he took for phlogiston, would turn back into metal. The post's charge fits the theory's later years, when its defenders adjusted it to each new finding.} Fake explanations do not feel fake; that is the danger.

``Modern research suggests,'' I say, that people reason about causes with something like the graphs of Bayes nets. \nb{Such research existed, for example Gopnik and colleagues' 2004 account of children's causal learning. The post cites none.} Rain makes the sidewalk wet, and a wet sidewalk is slippery. If we already know the sidewalk is wet, learning that it is slippery tells us nothing more about rain. Phlogiston makes fire hot and bright, and that feels like an explanation too.

One of the main inspirations for Bayesian networks, I say, was the problem of evidence being counted twice as inference passes back and forth between cause and effect. Judea Pearl's example is soldiers counting their line: each passes back a count of those ahead and passes forward a count of those behind, and the two are added only at the end. In the same way, evidence from the wet sidewalk goes back to rain and on to everything except the sidewalk. ``We count each piece of evidence exactly once; no update message ever `bounces' back and forth.'' \nb{That holds for networks without loops, such as the rain chain. With loops, Pearl wrote, messages ``may circulate indefinitely.''}

Phlogiston theorists let the evidence from hot fire go back to phlogiston and then counted it again as a prediction. Humans do not keep the two messages apart. A properly designed AI, I say, ``would notice the problem instantly,'' with ``just correct bookkeeping.'' \nb{That bookkeeping keeps one observation from being counted twice, which covers the bouncing. The post's other fault, a link set after looking at the result, it does not catch: the bookkeeping keeps no record of how a link was set.} ``Hindsight bias,'' I explain, is the nontechnical name for failing to separate the messages.

The phlogiston theorists were not fools, and your own fake explanations will not be labeled. So it is not enough to check how well a theory fits facts you already know: ``You've got to predict for tomorrow, not yesterday.'' It is ``the only way.'' \nb{Old facts have counted as strong evidence when a theory was not fitted to them. Mercury's perihelion, known for decades, counted for general relativity at least as much as its new prediction of light bending. Yudkowsky later wrote that Einstein used Mercury only ``for verification of his answer produced by other means.''}
